\documentclass[12pt]{article}
\usepackage{amsmath}
\usepackage{booktabs}
\usepackage{hyperref}
\usepackage{geometry}
\usepackage{graphicx}   % include graphics normally
\usepackage{float}
\geometry{margin=1in}

\title{Big Data Assignment Report}
\author{Jai Shankar}
\date{\today}

\begin{document}
\maketitle

\section*{Introduction}
I, Jai Shankar, prepared this report for the Big Data assignment.  It covers
k-gram generation, min-hashing, locality sensitive hashing (LSH), and experiments using the
MovieLens 100k dataset.  All results presented come from code that I executed in the
accompanying Jupyter notebook.

\section{Question 1: K-Grams}
\begin{figure}[H]
\centering
\includegraphics[width=0.8\linewidth,height=0.3\textheight,keepaspectratio]{q1.png}
\includegraphics[width=0.8\linewidth,height=0.3\textheight,keepaspectratio]{q2.png}
\caption{Output for Question 1}
\end{figure}
\FloatBarrier

For the four documents (D1.txt through D4.txt) the following counts of unique k-grams were
obtained:
\begin{tabular}{lccc}
\toprule
Document & \#char-2grams & \#char-3grams & \#word-2grams \\
\midrule
D1.txt & 263 & 765 & 279 \\
D2.txt & 262 & 762 & 278 \\
D3.txt & 269 & 828 & 337 \\
D4.txt & 255 & 698 & 232 \\
\midrule
Total (union of all documents) & 330 & 1301 & 743 \\
\bottomrule
\end{tabular}

Jaccard similarities between every pair of documents were computed for each k-gram type.
The results (18 values) appear in Tables~\ref{tab:jacc-c2}, \ref{tab:jacc-c3}, and
\ref{tab:jacc-w2}.

\textbf{Jaccard similarity using character 2-grams}

\begin{center}
\begin{tabular}{lcc}
\toprule
Pair & Similarity \\
\midrule
D1--D2 & 0.98113 \\
D1--D3 & 0.81570 \\
D1--D4 & 0.64444 \\
D2--D3 & 0.80000 \\
D2--D4 & 0.64127 \\
D3--D4 & 0.65299 \\
\bottomrule
\end{tabular}
\end{center}

\textbf{Jaccard similarity using character 3-grams}

\begin{center}
\begin{tabular}{lcc}
\toprule
Pair & Similarity \\
\midrule
D1--D2 & 0.97798 \\
D1--D3 & 0.58036 \\
D1--D4 & 0.30508 \\
D2--D3 & 0.56805 \\
D2--D4 & 0.30590 \\
D3--D4 & 0.31212 \\
\bottomrule
\end{tabular}
\end{center}

\textbf{Jaccard similarity using word 2-grams}

\begin{center}
\begin{tabular}{lcc}
\toprule
Pair & Similarity \\
\midrule
D1--D2 & 0.94077 \\
D1--D3 & 0.18234 \\
D1--D4 & 0.03024 \\
D2--D3 & 0.17366 \\
D2--D4 & 0.03030 \\
D3--D4 & 0.01607 \\
\bottomrule
\end{tabular}
\end{center}

\section{Question 2: Min-Hashing}
\begin{figure}[H]
\centering
\includegraphics[width=0.8\linewidth,height=0.3\textheight,keepaspectratio]{2.1.png}
\includegraphics[width=0.8\linewidth,height=0.3\textheight,keepaspectratio]{q2.2.png}
\caption{Output for Question 2}
\end{figure}

\subsection{Part 2A}
The pair D1, D2 with 3-grams has true Jaccard similarity $0.97798$.  Five signature sizes were
tried (t hash functions, using a large universe $m=100{,}000$):

\textbf{Approximate Jaccard values for different signature lengths}

\begin{center}
\begin{tabular}{ccc}
\toprule
t & Approximate Jaccard & Error \\
\midrule
20  & 0.97800 & 0.00002 \\
60  & 0.98040 & 0.00242 \\
150 & 0.97750 & 0.00048 \\
300 & 0.97930 & 0.00132 \\
600 & 0.97810 & 0.00012 \\
\bottomrule
\end{tabular}
\end{center}

These five numbers answer the first subquestion.  The hash family used was of the form
$h(x)=(ax+b)\bmod p$, with random coefficients and a prime $p>m$.

\subsection{Part 2B}
An extended experiment compared accuracy and running time over all six document pairs for
various values of $t$.  The mean error fell below $0.02$ around $t=150$, which also exhibited
reasonable runtime.  Therefore,
\textbf{recommendation: use $t=150$} as a good tradeoff between accuracy and speed.
(Details are documented in the notebook; the run for $t=150$ took approximately 0.1--0.2\,ms
per pair.)

\section{Question 3: Locality Sensitive Hashing}
\begin{figure}[H]
\centering
\includegraphics[width=0.8\linewidth,height=0.3\textheight,keepaspectratio]{q3.1.png}
\includegraphics[width=0.8\linewidth,height=0.3\textheight,keepaspectratio]{q3.2.png}
\caption{Output for Question 3}
\end{figure}

\subsection{Part 3A}
With $t=160$ total hash functions and threshold $\tau=0.7$, all factorisations $b\times r=160$
were examined.  The configuration that maximised separation around $\tau$ was
\textbf{$b=10$ rows per band, $r=16$ bands}.  This choice gives a sharp S-curve and was used
for subsequent probability calculations.

\subsection{Part 3B}
Using the above $(b,r)$ the probability that each of the six document pairs would collide in
at least one band was computed via
$f(s)=1-(1-s^{b})^{r}$.  The results are:
\begin{tabular}{lll}
\toprule
Pair & Jaccard ($s$) & Collision prob. \\
\midrule
D1--D2 & 0.9780 & 0.99  \\
D1--D3 & 0.5804 & 0.04  \\
D1--D4 & 0.3051 & 0.00  \\
D2--D3 & 0.5680 & 0.03  \\
D2--D4 & 0.3059 & 0.00  \\
D3--D4 & 0.3121 & 0.00  \\
\bottomrule
\end{tabular}

These six probabilities answer the final subquestion of Part~3.

\textbf{LSH collision probabilities for document pairs}

\begin{center}
\begin{tabular}{lll}
\toprule
Pair & Jaccard ($s$) & Collision prob. \\
\midrule
D1--D2 & 0.9780 & 0.99  \\
D1--D3 & 0.5804 & 0.04  \\
D1--D4 & 0.3051 & 0.00  \\
D2--D3 & 0.5680 & 0.03  \\
D2--D4 & 0.3059 & 0.00  \\
D3--D4 & 0.3121 & 0.00  \\
\bottomrule
\end{tabular}
\end{center}

These six probabilities answer the final subquestion of Part~3.

\section{Question 4: MovieLens Experiments}
\begin{figure}[htbp]
\centering
\includegraphics[width=0.8\linewidth]{q4.png}
\caption{Output for Question 4}
\end{figure}

The MovieLens 100k data file was downloaded from the official URL and processed.  There

\section*{Repository}
The complete notebook, data files, and supplementary material for this assignment are hosted
on GitHub at \url{https://github.com/jaishankar02/m25csa014_ML_with_Bigdata_Assgn2}.

The MovieLens 100k data file was downloaded from the official URL and processed.  There
are 943 users and 1682 distinct movies.

Exact Jaccard similarities were computed for every user pair; 10 pairs have similarity at
least 0.5.  A sample of the top pairs is listed in the notebook.

\textbf{MovieLens min-hash experiment averages}

\begin{center}
\begin{tabular}{ccccc}
\toprule
k & Avg.~estimated pairs & Avg.~false pos. & Avg.~false neg. \\
\midrule
50  & 131.2 & 123.0 & 1.8 \\
100 & 32.2 & 28.0  & 0.8 \\
200 & 16.4 & 17.0  & 0.4 \\
\bottomrule
\end{tabular}
\end{center}

The approximate counts approach the exact value as $k$ increases.

\section{Question 5: LSH on MovieLens}
\begin{figure}[htbp]
\centering
\includegraphics[width=0.8\linewidth]{q5.png}
\caption{Output for Question 5}
\end{figure}

Using the configurations specified in the problem statement and thresholds 0.6 and 0.8, the
number of false positives and false negatives (averaged over five runs) were recorded.

For threshold $\ge0.6$ there were 3 true pairs; for $\ge0.8$ only 1 true pair.

Results:

\textbf{LSH false positives / negatives (averaged)}

\begin{center}
\begin{tabular}{lllll}
\toprule
Threshold & k & (r,b) & Avg.~FP & Avg.~FN \\
\midrule
0.6 & 50  & (5,10) & 0.8 & 2.2 \\
0.6 &100  & (5,20) & 0   & 3.0 \\
0.6 &200  & (5,40) & 0   & 3.0 \\
0.6 &200  & (10,20)& 0   & 3.0 \\
0.8 & 50  & (5,10) & 1.0 & 0.4 \\
0.8 &100  & (5,20) & 0   & 1.0 \\
0.8 &200  & (5,40) & 0   & 1.0 \\
0.8 &200  & (10,20)& 0   & 1.0 \\
\bottomrule
\end{tabular}
\end{center}

Only the small signature table ($k=50$) produced any candidate pairs; larger tables were too
strict with the chosen banding and yielded zero candidates.

\section*{Conclusion}
This report documents the implementation and results of k-gram extraction, min-hashing,
LSH, and analysis on both the supplied documents and the MovieLens dataset.  All code and
experiments were executed by me, Jai Shankar, in the Jupyter notebook accompanying this
write-up.

\end{document}
